\subsection{对数曲线}%
\label{sub:Log curves}

目前，所有幺半群和对数概形都取为 fine 的（甚至是 fine saturated 的）。理由有二：
\begin{enumerate}
    \item 若 \(P\) 是 fine 幺半群，则其对偶是 fs 且尖（sharp）的；
    \item 若 \(X = (\ul{X}, M_X)\) 是 fine 对数概形，则对任意几何点 \(\bar{x} \to X\)，存在 \(\bar{x}\) 的一个 \'etale 邻域带有图表。
\end{enumerate}

下面讨论幺半群之间映射的各种几何性质。
\begin{defn}
    称映射 \(\theta \colon P \to Q\) 是 \textit{integral}（对数积）的，若映射 \(\Spec\C[Q] \to \Spec \C[P]\) 是平坦的。
\end{defn}

例如，任意非零映射 \(\N \to \N\) 以及对角映射 \(\Delta \colon\N \to \N^2\) 都是 integral 的。若一个映射是 integral 的，则对任意 \(P \to P'\)，幺半群 \(P' \oplus_P Q\) 是 fine 的。

\begin{exm}
    推出未必是 fine 的，例如考虑
    \[ \ab< a ,b,c \mid a+c=b+c> = \N^3 \oplus_{\N^2} \N. \]
\end{exm}

\begin{defn}
    fine 对数概形之间的映射 \(f \colon X \to Y\) 称为 integral 的，若在 \(X\) 的每个几何点处，映射
    \[ (f^*\bar{M}_Y)_{\bar{x}} \to \bar{M}_{X,\bar{x}} \]
    是 integral 的。
\end{defn}

\begin{defn}
    fs 幺半群之间的映射 \(P \to Q\) 称为 \textit{saturated}（饱和）的，若对所有 \(P \to P'\)，幺半群 \(P' \oplus_P Q\) 也是 fs 的。
\end{defn}

\begin{exm}
    映射 \(\Delta \colon \N \to \N^2\) 是 saturated 的，但 \(\ell \colon \N \to \N\) 不是。
\end{exm}

\begin{defn}
    fs 对数概形 \(S\) 上的 \textit{对数曲线} 是指 fs 对数概形之间的对数光滑且 integral 的映射 \(\pi \colon \mc{C} \to S\)，使得其底几何纤维是既约连通曲线。
\end{defn}

\begin{prop}[{\cite{Kato1989}}]
    设 \(f \colon X \to Y\) 是 fs 对数概形之间的对数光滑且 integral 的映射。则 \(\ul{f}\) 是平坦的。
\end{prop}

\begin{thm}[{\cite{FKato2000}}]
    对数曲线的底族是结点曲线族。
\end{thm}

\begin{exm}
    设 \(P = \N \setminus \ab\{1\}\)。则 \(\Spec (P \to k[P])\) 在 \(k\) 上对数光滑，但 \(\Spec k[P]\) 是一个尖点。
\end{exm}

\begin{defn}
    fs 基 \(S\) 上的 \textit{\(n\)-点对数曲线} 由数据 \((\pi \colon \mc{C} \to S), \vec{p} = \ab\{p_i\}_{i=1}^n\) 组成，满足
    \begin{enumerate}
        \item 底族是预稳定曲线；
        \item \(\pi\) 是 fs 对数概形之间的对数光滑映射；
        \item 在 \(\ul{\pi}\) 的光滑轨迹 \(U\) 上，有
            \[ \bar{M}_{\mc{C}}|_U \simeq \pi^* \bar{M}_S \oplus \bigoplus_{i=1}^n p_i \ul{\N}. \]
    \end{enumerate}
\end{defn}

\subsection{预稳定曲线上的典范对数结构}%
\label{sub:Canonical log structures on prestable curves}

设 \(\mf{M}_{g,n}\) 为 \(n\) 点预稳定曲线的模空间。则奇异轨迹是一个正规交叉除子，故有典范对数结构
\[ M_{\mf{M}_{g,n}}^{\on{can}} = M_{\partial \mf{M}_{g,n}}. \]
若考虑泛族 \(\mf{C}_{g,n}\)，可以构造 \(\pi^{-1}(\partial \mf{M}_{g,n}) \subseteq \mf{C}_{g,n}\)。于是我们有映射
\[ \pi^* M_{\partial \mf{M}_{g,n}} \to M_{\pi^{-1}(\partial \mf{M}_{g,n})}. \]
也可以考虑除子 \(p_i(\mf{M}_{g,n}) \subset \mf{C}_{g,n}\)，它们同样定义对数结构
\[ P_i \coloneqq M_{p_i(\mf{M}_{g,n})}. \]
因此可以定义对数结构
\[ M_{\mf{C}_{g,n}}^{\on{can}} \coloneqq M_{\pi^{-1}(\partial\mf{M}_{g,n})} \oplus_{\mc{O}^{\times}} \bigoplus_{i=1}^n P_i. \]
于是若考虑
\[ \pi^{\flat} \colon \ul{\pi}^{-1} M_{\mf{M}_{g,n}}^{\on{can}} \to M_{\pi^{-1}(\partial \mf{M}_{g,n})} \to M_{\mf{C}_{g,n}}^{\on{can}}, \]
则可将 \(\pi\) 提升为对数叠之间的映射。

\begin{thm}[{\cite{FKato2000}}]
    映射 \(\pi\) 是对数光滑的 integral 映射。
\end{thm}

\begin{defn}
    对于预稳定曲线族 \((\ul{C}/\ul{S}, \vec{p})\)，可沿如下图表拉回
    \begin{equation*}
    \begin{tikzcd}
        \ul{C} \ar{r} \ar{d} & \mf{C}_{g,n} \ar{d} \\
        S \ar{r} & \mf{M}_{g,n}
    \end{tikzcd}
    \end{equation*}
    并在左侧的边上赋予典范对数结构。
\end{defn}

\begin{rmk}
    设 \(\ul{S} = \Spec k\)。若 \(S^{\on{can}} = \Spec (\N^r \to k)\)，则 \(r\) 是 \(\ul{C}\) 上结点的个数。在结点 \(q\) 处 \'etale 局部地看，对数结构形如
    \begin{equation*}
    \begin{tikzcd}
        \N^{r-1} \oplus \N^2 \ar{r} & k[x,y]/xy \\
        \N^{r-1} \oplus \N \ar{u}{\on{id}\oplus\Delta} \ar{r} & k \ar{u}.
    \end{tikzcd}
    \end{equation*}
\end{rmk}

\begin{thm}[{\cite{FKato2000}}]
    设 \((\mc{C}/S, \vec{p})\) 是 fs 基 \(S\) 上的 \(n\) 点对数曲线。则存在唯一的 \(M_S^{\on{can}} \to M_S\) 使得
    \begin{equation*}
    \begin{tikzcd}
        \mc{C} \ar{d} \ar{r} & \mc{C}^{\on{can}} \ar{d} \\
        S \ar{r}{a} & S^{\on{can}}
    \end{tikzcd}
    \end{equation*}
    是 fs 对数概形范畴中的笛卡尔方图。
\end{thm}

一句口号是：对数曲线就是预稳定曲线连同 \(M_S^{\on{can}} \to M_S\)。

\begin{exm}
    设 \(\ul{S} = \Spec k[s]\)，\(\ul{T} = \Spec k[t]\)。再取 \(\ell \in \N\)。考虑图表
    \begin{equation*}
    \begin{tikzcd}
        \Spec \frac{k{ [x,y,s] }}{xy-s^{\ell}} \ar{d} \ar{r} & \Spec \frac{k{ [x,y,s] }}{xy-s^{\ell}} \ar{d} \ar{r} & \Spec  \frac{k{[x,y,t]}}{xy-t} \ar{r}\ar{d} & \mf{C}_{g,n} \ar{d} \\
        \Spec k{[s]} \ar{r}{\ul{ \on{id} }} & \Spec {k[s]} \ar{r} & \Spec k{[t]} \ar{r} & \mf{M}_{g,n}.
    \end{tikzcd}
    \end{equation*}
    在对数结构上，考虑左下角箭头对应 \(\N\) 上 \(1 \mapsto \ell\) 的情形。此时左上角的幺半群变为 \(\N^2 \oplus_{\N, \ell} \N\)。记之为 \(P\)，则映射 \(k[P] \to k[\N]\) 不是带典范对数结构的对数曲线族，但确实是对数曲线族。
\end{exm}

\subsection{对数曲线的模空间}%
\label{sub:Moduli of log curves}

\begin{defn}
    \(n\) 点对数曲线之间的映射 \((\mc{C}, S) \to (\mc{C}', S')\) 是指一对 \((\ell \colon \mc{C} \to \mc{C}', \Theta \colon S \to S')\)，其中 \(\Theta\) 是严格的，且图表
    \begin{equation*}
    \begin{tikzcd}
        \mc{C} \ar{d} \ar{r}{\ell} & \mc{C}' \ar{d} \\
        S \ar{r}{\Theta} & S'
    \end{tikzcd}
    \end{equation*}
    是笛卡尔的。记 \(\mf{M}_{g,n}^{\log}\) 为参数化 \(n\) 点对数曲线的纤维范畴。
\end{defn}

\begin{prop}[{\cite{FKato2000,Olsson2003}}]
    我们有 \(\mf{M}_{g,n}^{\log} \simeq \ms{Log}_{(\mf{M}_{g,n}, M_{\mf{M}_{g,n}}^{\on{can}})}\)。
\end{prop}

特别地，对概形 \(\ul{T}\)，群胚 \(\mf{M}_{g,n}^{\log}(\ul{T})\) 参数化对数结构 \(M_T\) 连同对数曲线 \(\mc{C} / T\)。

考虑 \(S = \Spec (Q \to k)\) 上的对数曲线 \(\pi \colon \mc{C} \to S, \vec{p}\)。记 \(\bar{\pi}^{\flat} \colon \bar{\pi}^* \bar{M}_S \to \bar{M}_{\mc{C}}\)。在一般点或一般的光滑点 \(\eta\) 处，映射为
\[ \bar{\pi}_{\eta}^{\flat} \colon (\pi^{-1} \bar{M}_S)_{\eta} = Q \simeq \bar{M}_{\mc{C}, \eta}. \]
在标记点 \(p_i \in \ol{\ab\{\eta\}}\) 处，有 \(\bar{M}_{\mc{C}, p_i} = Q \oplus \N\)。此时有
\[ \chi_{\eta} \colon \bar{M}_{\mc{C}, p_i} \xrightarrow{\on{pr}_1} \bar{M}_{\mc{C}, \eta} \]
这一特化映射，以及拉回
\[ \bar{\pi}_{p_i}^{\flat} \colon Q \hookrightarrow Q \oplus \N. \]
在结点 \(q\) 处，有 \(\bar{M}_{\mc{C}, q} = Q \oplus_{\N} \N^2\)。记 \(\rho_q\) 为 \(1\) 在 \(Q\) 中的像，它是 \(q\) 处的边长或光滑化参数。再记 \(\bar{\eta}_1\) 与 \(\bar{\eta}_2\) 为 \(q\) 两个分支的一般点。则一般化映射为
\begin{align*}
    \chi_{\eta_1} \colon Q \oplus_{\N} \N^2 \to Q \qquad (q, (a,b)) \mapsto q + a \rho_q ;\\
    \chi_{\eta_2} \colon Q \oplus_{\N} \N^2 \to Q \qquad (q, (a,b)) \mapsto q + b \rho_q .
\end{align*}

\subsection{对数曲线的热带化}%
\label{sub:Tropicalization of log curves}

我们只考虑对数点 \(S = \Spec(Q \to k)\) 上的情形。考虑锥范畴 \(\ms{Cone}\)，其对象是满足 \(w \subseteq N_w \otimes_{\Z} \R\) 的对 \((w, N_w)\)，其中 \(w\) 是满维强凸有理多面锥。态射是映射 \((w_1, N_{w_1}) \to (w_2, N_{w_2})\)，由
\[ \chi \colon N_{w_1} \to N_{w_2} \]
给出，满足 \(\chi_{\R}(w_1) \subseteq w_2\)。

\begin{exm}
    设 \(X\) 是 fs 对数概形，考虑 \(( (\bar{M}_{X,x}^{\vee})_{\R}, (\bar{M}_{X,x}^{\vee})_{\Z} )\)。
\end{exm}

现在考虑广义锥复形的范畴 \(\ms{GCC}\)。其对象是拓扑空间 \(\abs{\Sigma}\) 连同它作为 \(\ms{Cone}\) 中一个图表之余极限的表示，其中所有箭头都是面的包含。其态射是连续映射 \(\abs{\Sigma} \to \abs{\Sigma'}\)，使得对每个 \((w, N_w) \in \Sigma\)，映射 \(w \to \abs{\Sigma'}\) 通过某个 \((w, N_w) \to (w', N_{w'})\) 分解。

\begin{exm}
    若 \(f \colon X \to Y\) 是 fs 对数概形的态射，则映射 \(\bar{f}^{\flat} \colon f^* \bar{M}_Y \to \bar{M}_X\) 在热带化上诱导 \(\Sigma(X) \to \Sigma(Y)\)。
\end{exm}

\begin{defn}
    称广义锥复形 \(\Sigma\) 为 \textit{锥复形}，若对每个 \((\sigma, N_{\sigma}) \in \Sigma\)，映射 \(\sigma \to \abs{\Sigma}\) 是单的。
\end{defn}

\begin{thm}[{\cite{ACP2015}}]
    若 \(X\) 是 fs 的，且在带 Zariski 对数结构的 \(\ul{\C}\) 上对数光滑，则 \(\Sigma(X)\) 是锥复形。
\end{thm}

对幺半群 \(Q\)，记 \(w_Q = Q_{\R}^{\vee} = \Hom(Q, \R_{\geq 0})\)，格为 \(\Hom(Q, \Z) = Q_{\Z}^{\vee}\)。在光滑点处，有
\[ w_{\eta} = (Q_{\R}^{\vee}, Q_{\Z}^{\vee}). \]
在标记点处，可见
\[ w_p = \ab(\Hom (Q \oplus \N, \R_{\geq 0}) = Q_{\R}^{\vee} \oplus \R_{\geq 0}, Q_{\Z}^{\vee}\oplus \Z). \]
在结点处，有
\begin{align*}
    w_q &= \ab(\Hom(Q \oplus_{\N} \N^2, \R_{\geq 0}), ( \bar{M}_{\mc{C}, q} )_{\Z}^{\vee}).
\end{align*}
注意该锥实际上等于
\[ \ab\{(m, m_x, m_y) \mid m(\rho_q) = m_x + m_y \} \simeq \ab\{(m, e_x) \mid e_x \leq m(\rho_q) \} \subseteq Q_{\R}^{\vee} \times \R_{\geq 0}. \]

既然已经弄清了与每个点相伴的锥，接下来需要讨论面映射。与标记点相对应的一般化映射为
\[ Q \oplus \N \xrightarrow{\on{pr}_1} Q, \]
相应的面映射为
\[ w_{\eta} \hookrightarrow w_p \qquad m \mapsto (m,0). \]
在结点处，两个面映射为
\[ \chi_{\eta_1}^{\vee} \colon Q_{\R}^{\vee} \to w_q \qquad m \mapsto (m, m(\rho_q)) \]
与
\[ \chi_{\eta_2}^{\vee} \colon Q_{\R}^{\vee} \to w_q \qquad m \mapsto (m, 0). \]

\begin{exm}
    考虑一条由两条有理曲线并成的曲线，在其中一支上有一个标记点 \(p\)。取 \(\ul{S} = \Spec \C\)。于是可见
    \[ S^{\on{can}} = \Spec(\N \to \C), \]
    \(Q = \N\)，且 \(Q_{\R}^{\vee} = \R_{\geq 0}\)。还可见 \(\bar{M}_{C, q} = \N \oplus_{\N} \N^2\) 以及 \(\bar{M}_{C, p} = \N \oplus \N\)。
\end{exm}

\begin{figure}[H]
    \centering
    \begin{tikzpicture}[
        x=1cm,
        y=1cm,
        cone ray/.style={thick,->,>=stealth},
        common face/.style={red,line width=2pt,line cap=round},
        diagonal face/.style={blue!75!black,line width=2pt,line cap=round}
    ]
        % The cone at the marked point.
        \coordinate (p) at (0,4.45);
        \draw[cone ray] (p) -- ++(-0.55,2.15);
        \draw[cone ray] (p) -- ++(2.25,1.85);
        \draw[diagonal face]
            ($(p)+(0.75,0.62)$) -- ($(p)+(1.45,1.19)$);
        \node[above] at ($(p)+(0.55,2.2)$)
            {$w_p=\R_{\geq 0}\times\R_{\geq 0}$};

        % The cone at the first generic point.
        \coordinate (etaone) at (0.35,4);
        \draw[cone ray] (etaone) -- ++(2.25,1.85);
        \draw[diagonal face]
            ($(etaone)+(0.75,0.62)$) -- ($(etaone)+(1.45,1.19)$);
        \node[above right] at ($(etaone)+(2.0,1.6)$) {$w_{\eta_1}$};

        % The cone at the node.
        \coordinate (q) at (0.35,3.25);
        \draw[cone ray] (q) -- ++(2.25,1.85);
        \draw[cone ray] (q) -- ++(3.35,0);
        \draw[diagonal face]
            ($(q)+(0.75,0.62)$) -- ($(q)+(1.45,1.19)$);
        \draw[common face] ($(q)+(0.85,0)$) -- ($(q)+(1.75,0)$);
        \node at ($(q)+(2.55,0.75)$) {$w_q$};
        \node[anchor=west] at ($(q)+(4.05,0.75)$)
            {$\bigl\{(a,b)\mid b\leq a\bigr\}\subseteq\R_{\geq 0}^{2}$};

        % The colored intervals indicate the common faces used in the gluing.
        \coordinate (etatwo) at (0.35,2.75);
        \draw[cone ray] (etatwo) -- ++(3.35,0);
        \draw[common face] ($(etatwo)+(0.85,0)$) -- ($(etatwo)+(1.75,0)$);
        \node[right] at ($(etatwo)+(3.45,0)$) {$w_{\eta_2}$};
    \end{tikzpicture}
\end{figure}


\subsection{热带曲线}%
\label{sub:Tropical curves}

\begin{defn}
    锥 \(w, N_w\) 上的 \textit{热带曲线} 族 \((G, \vec{g}, \vec{\ell})\) 由如下数据组成
    \begin{enumerate}
        \item 连通图 \(G = V(G) \cup E(G) \cup L(G)\)；
        \item 腿的排序 \(\vec{m} \colon L(G) \xrightarrow{\sim} \ab\{1, \ldots, n\}\)；
        \item 亏格函数 \(\vec{g} \colon V(G) \to \N\)；
        \item 边长函数 \(\ell \colon E(G) \to \Hom(w \cap N_w, \N) \setminus \ab\{0\}\)。
    \end{enumerate}
\end{defn}

现在构造 \(\Sigma(G, \ell)\)。对 \(v \in V(G)\)，置 \(w_v = w\)。对 \(e \in E(G)\)，置
\[ w_e = \ab\{ (s, \lambda) \mid \lambda \leq \vec{\ell}(e)(s) \} \subseteq w \times \R_{\geq 0}. \]
对腿 \(\ell \in L(G)\)，置 \(w_{\ell} = w \times \R_{\geq 0}\)。现在需要定义面映射。若 \(\ell\) 与 \(v\) 关联，则面映射为
\[ w_v \hookrightarrow w_{\ell} \colon s \mapsto (s, 0). \]
若 \(e\) 与两个顶点 \(v_1, v_2\) 关联，则面映射为
\begin{align*}
    w_{v_1} \hookrightarrow w_e &\colon s \mapsto (s, \vec{\ell}(e)(s)) \\
    w_{v_2} \hookrightarrow w_e & \colon s \mapsto (s, 0).
\end{align*}
\begin{rmk}
    \(v_1\) 与 \(v_2\) 的选取不改变锥复形（在同构意义下），但在为对数映射定义接触阶时它确实有影响。然而一般而言，没有典范的方式来作此选取。
\end{rmk}

\subsection{穿孔曲线}%
\label{sub:Punctured curves}

穿孔对数曲线与穿孔映射发展于~\cite{ACGS2025}。

考虑结点曲线
\[ \Spec k[x,y]/xy \to \Spec k \]
其对数结构由 \(\N^2\) 给出，图表为 \((a,b) \mapsto x^a y^b\)，而 \(\N \to k\) 是标准对数点。现在考虑 \(\ul{C}_1 = Z(x) = \Spec k[y]\)，并以 \(p \colon \Spec k \to \Spec k[y]\) 标记该结点。注意
\[ C^{\on{can}} = (\ul{C}, M_C) \to S^{\on{can}} = \Spec(\N \to k). \]
记 \(M_{C_1}^{\circ} = M_C|_{C_1}\)。则在对数结构上有图表
\begin{equation*}
\begin{tikzcd}
    k[y] & & \N^2 \ar{ll}[swap]{(a,b) \mapsto 0^a y^b} \\
    k \ar{u} & & \N \ar{ll} \ar{u}{\Delta}.
\end{tikzcd}
\end{equation*}
特别地，可以计算茎 \(\bar{M}_{C_1}^{\circ}|_p = \N^2\)，因此有
\[ \bar{\pi}_{C_1^{\circ}}^{\flat} |_p \colon \N \xrightarrow{\Delta} \N^2. \]
然而本应有映射
\[ \bar{\pi}^{\flat}_{C_1}|_p \colon \N \xrightarrow{i_1} \N \oplus \N, \]
故 \(C_1\) \textbf{不是} 对数曲线。于是可以考虑图表
\begin{equation*}
\begin{tikzcd}
    \N \oplus \Z \\
    \N^2 \ar[hookrightarrow]{u} \ar{rr}{(0,1)\mapsto (0,1)}[swap]{(1,0)\mapsto (1,1)} & & \N^2 \ar{ull}[swap]{\Vectorstack{{(1,0)\mapsto(1,-1)} {(0,1)\mapsto(0,1)}}} \\
    & \N \ar{ul}{i_1} \ar{ur}{\Delta}.
\end{tikzcd}
\end{equation*}
因此我们将考虑扩张对数结构。

回顾对数曲线满足 \(M_{\mc{C}} = M_{\mc{C}}' \oplus_{\mc{O}_{\mc{C}}^{\times}} P_{\mc{C}}\)，其中 \(P_{\mc{C}}\) 是由标记给出的除子对数结构。

\begin{defn}
    fs 对数概形 \(S\) 上的 \textit{穿孔曲线} 由 \((\mc{C}^{\circ} \xrightarrow{p} \mc{C} \xrightarrow{\pi} S, \{p_i\}_{i=1}^n)\) 组成，其中 \((\mc{C} \xrightarrow{\pi} S, \vec{p})\) 是对数曲线，且映射 \(p \colon \mc{C}^{\circ} \to \mc{C}\) 是 fine 对数概形之间的映射，满足下列条件：
    \begin{enumerate}
        \item 底映射 \(\ul{p}\) 是同构；
        \item 在对数结构上，映射 \(p^{\flat} \colon M_{\mc{C}} \hookrightarrow M_{\mc{C}^{\circ}}\) 在 \(p_i\) 之外是同构，并诱导
            \[ M_{\mc{C}} \xrightarrow{p^{\flat}} M_{\mc{C}^{\circ}} \hookrightarrow M_{\mc{C}} \oplus P^{\mr{gp}}. \]
        \item 对任意标记 \(p_i\) 及任意 \(s \in M_{\mc{C}^{\circ}, p_i} \setminus M_{\mc{C}, p_i}\)，有 \(\alpha_{\mc{C}^{\circ}}(s) = 0\)。
    \end{enumerate}
\end{defn}

粗略地说，穿孔曲线就是对数曲线加上一个交换图表
\begin{equation*}
\begin{tikzcd}
    M_{\mc{C}^{\circ}} \ar{dr} \ar[hookrightarrow]{r} & M_{\mc{C}}' \oplus_{\mc{O}_{\mc{C}}^{\times}} P^{\on{gp}} \\
    & \mc{O}_{\mc{C}} \\
    M_{\mc{C}} \ar[hookrightarrow]{uu} \ar{ur}.
\end{tikzcd}
\end{equation*}


\begin{exm}
    设 \((\ul{\mc{C}} \to \ul{S} = \Spec k, p)\) 是光滑的 \(1\)-点曲线。
    \begin{enumerate}
        \item 设 \((\mc{C}/S, p)\) 赋予典范对数结构（即 \(S\) 无对数结构）。令 \(M_{\mc{C}} = P\)。则任何穿孔 \(p^{\flat} \colon M_{\mc{C}} \hookrightarrow M_{\mc{C}^{\circ}} \subset P^{\on{gp}}\) 必为平凡。原因在于 \(M_{\mc{C}^{\circ}}\) 不能有截面 \(s\) 使得 \(\alpha_{\mc{C}^{\circ}}(s) = 0\)。若 \(s \in P = M_{\mc{C}}\)，则 \(\alpha_{\mc{C}^{\circ}}(s) = \alpha_{\mc{C}}(s)\)。否则若 \(s \notin P\)，则因 \(\alpha_{\mc{C}^{\circ}}(-s)\) 非零，\(\alpha_{\mc{C}^{\circ}}(s)\) 也非零。
        \item 设 \(M_{\mc{C}^{\circ}} \subseteq (\N \oplus \mc{O}_{\mc{C}}^{\times}) \oplus P^{\on{gp}}\) 是一个穿孔。取 \(\mc{O}(-p)\) 的一个局部生成元 \(x\)。则任意 \(s \in M_{\mc{C}^{\circ}, p}\) 具有形式 \(s = ((n, \varphi), x^m)\)，其中 \(n \in \N\)，\(\varphi \in \mc{O}_{\mc{C}, p}^{\times}\)。若 \(m < 0\)，则我们要求 \(\alpha_{\mc{C}^{\circ}}(s) = 0\)。这意味着 \(\alpha_{\mc{C}}((n, \varphi), 0) = 0\)，因此
            \[ \bar{M}_{\mc{C}^{\circ}, p} \subseteq \ab\{(n, m) \in \N \oplus \Z \mid m \geq 0 \text{ if } n = 0 \}. \]
            事实上，右端任意包含 \(\N \oplus \N\) 的 fine 子幺半群均可由某个穿孔实现。
        \item 设 \(S = \Spec(\N \to k)\)。考虑 \(M_{\mc{C}^{\circ}}\) 使得 \(\bar{M}_{\mc{C}^{\circ}, p} = \N^2 + \N (2,-2)\)。由于 \((1,-1) \notin \bar{M}_{\mc{C}^{\circ}, p}\)，可见 \(M_{\mc{C}^{\circ}}\) 不是饱和的。
    \end{enumerate}
\end{exm}

\begin{lem}[{\cite{ACGS2025}}]
    设 \(Q\) 是尖的 fs 幺半群。令 \(w \coloneqq \Hom(Q, \R_{\geq 0}) = Q_{\R}^{\vee}\)。设 \(Q^{\circ} \subseteq Q \oplus \Z\) 有限生成，且 \(Q \oplus \N \subsetneq Q^{\circ}\) 以及 \(Q^{\circ} \cap \ab\{0\} \times \Z_{<0} = \emptyset\)。则存在非零的凹的分段线性函数
    \[ \ell \colon w \to \R_{\geq 0} \]
    具有有理斜率，使得
    \[ (Q^{\circ})_{\R}^{\vee} = \ab\{ (s, \lambda) \in w \times \R_{\geq 0} \mid 0 \leq \lambda \leq \ell(s) \}. \]
\end{lem}

\begin{exm}
    设 \(Q = \N^2\)，置
    \begin{align*}
        Q^{\circ} &= \ab<((1,0), -1), ((0,1), -1), ((0,0),1)> \subset \N^2 \oplus \Z \\
        &= \ab\{((a,b), m) \mid a+b+m \geq 0, a,b \geq 0 \}.
    \end{align*}
    则
    \[ (Q^{\circ})_{\R}^{\vee} = \ab\{ (e_x, e_y, e_t) \mid 0 \leq e_t \leq \min(e_x, e_y) \}. \]
\end{exm}





